\documentclass[a4paper,12pt]{article}
\usepackage[utf8]{inputenc}
\usepackage[spanish]{babel}
\usepackage{color}
\usepackage{parskip}
\usepackage{graphicx}
\usepackage{multirow}
\usepackage{listings}
\usepackage{vmargin}
\usepackage{datetime}
\newdate{date}{24}{09}{2017}
\graphicspath{ {imagenes/} }
\definecolor{mygreen}{rgb}{0,0.6,0}
\definecolor{lbcolor}{rgb}{0.9,0.9,0.9}
\usepackage{epstopdf}


\lstset{
backgroundcolor=\color{lbcolor},
    tabsize=4,    
%   rulecolor=,
    language=[GNU]C++,
        basicstyle=\tiny,
        aboveskip={1.5\baselineskip},
        columns=fixed,
        showstringspaces=false,
        extendedchars=false,
        breaklines=true,
        prebreak = \raisebox{0ex}[0ex][0ex]{\ensuremath{\hookleftarrow}},
        frame=single,
        showtabs=false,
        showspaces=false,
        showstringspaces=false,
        identifierstyle=\ttfamily,
        keywordstyle=\color[rgb]{0,0,1},
        commentstyle=\color[rgb]{0.026,0.112,0.095},
        stringstyle=\color{red},
        numberstyle=\color[rgb]{0.205, 0.142, 0.73},
%        \lstdefinestyle{C++}{language=C++,style=numbers}’.
}


\begin{document}
\title{Pursuit Internet}
\author{
Christofer Fabián Chávez Carazas \\
\small{Universidad Nacional de San Agustín de Arequipa} \\
\small{Escuela Profesional de Ciencia de la Computación} \\
\small{Computación Centrada en Redes}
}
\date{\displaydate{date}}

\maketitle

Desde los inicios de la Internet, ésta a funcionado siempre bajo la misma arquitectura: ordenadores que se conectan a servidores
desde donde se descargan y se transmiten los datos. Esto a provocado que surjan un sin fin de problemas, tales como: SPAM,
ataques de denegación de servicio y poco desempeño a la hora de enviar datos multimedia \cite{pursuitPDF}. Varios investigadores
han trabajado en la tarea de crear una nueva arquitectura para Internet que transforme la forma en que la información es compartida,
y que se haga de forma rápida y segura, evitando los problemas de la arquitectura cliente-servidor. \par
El prototipo, financiado por la Unión Europea y desarrollado por un grupo de investigadores de la Universidad de Cambridge, se
llama PURSUIT, un nuevo protocolo basado en el P2P. En esta nueva arquitectura, en lugar de conectarse a servidores centrales y organizar toda la arquitectura en torno a ellos,
como ocurre en la actualidad, PURSUIT propone descentralizarla y lograr que los ordenadores y equipos individuales sean capaces de recibir, copiar y re-publicar contenido que reciben de otros equipos.
Es decir, los ordenadores podría acceder a datos e información de múltiples fuentes a la vez, en lugar de un solo servidor central. \cite{UC}.
Al tener los datos en una gran número de ordenadores, se garantiza que no se pueda eliminar su acceso tan fácilmente, y cuantas más personas accediesen a esos archivos, más fácilmente sería accesible y más complicada su eliminación.
La idea es parecida a la que usan los clientes de torrent. Cuando descargamos un archivo por torrent, no lo estamos obteniendo de un único lugar,
sino que conseguimos varios trozos inconexos dispersos por toda la red, en los discos duros de todos los usuarios que estén usando ese torrent. Una vez que tenemos los trozos,
el programa los junta formando el archivo que estábamos buscando \cite{Omi}. En la Figura \ref{fig:PvsI} se muestra una comparación
entre la Internet actual y Pursuit internet.


\begin{figure}
 \centering
 \includegraphics[scale = 0.5]{1.jpg}
 \caption{Pursuit internet vs Internet \cite{UC}}
 \label{fig:PvsI}
\end{figure}

Los investigadores proponen lanzar esta arquitectura a escala global, de forma que la anterior estructura conectada a servidores,
podría sustituirse por otra de equipos conectados entre sí. En definitiva, se trata de eliminar el concepto de la URL de Internet.
Su propuesta es que Internet deje de funcionar en torno a URL y pase a funcionar en base a URI (\textit{Uniform Resource Identifiers}).
Estas URIs especificarían dónde se encuentran los datos y qué ruta seguir en la nueva red para encontrarlos y descargarlos \cite{Omi}. \\
Los investigadores aseguran que esta nueva arquitectura sería mucho más segura porque tú podrías controlar quién tiene acceso
a la información personal que tú envías. Sin embargo, también hay múltiples obstáculos, entre ellos, el aún escaso poder de
computación de móviles y tabletas comparado con los ordenadores, o el enorme reto tecnológico que supondría pasar de una
arquitectura a otra. \par

PURSUIT se basa y continúa con el trabajo hecho en el proyecto PSIRP, que utiliza el paradigma publicar/subscribir (pub/sub) \cite{ps}.
Pub/sub es un paradigma centrado en la información que cambia el poder lejos del emisor de datos, es decir, los consumidores de datos expresan su inter´es
en información específica, que les envía la red cuando se vuelve disponible. Como resultado, la información se propaga sólo a los nodos, que realmente quieren acceder
a él. En una red pub/sub, el emisor ``publica'' los datos que quiere enviar, y el receptor se ``subscribe'' a las publicaciones
que quiere recibir. Esto quiere decir que no se recibe ningún dato al cual no se haya suscrito con anterioridad.

\newpage

\begin{large}
 \textbf{Conclusiones}
\end{large}

\begin{itemize}
 \item PURSUIT es un prototipo que podría ayudar a resolver varios problemas en la arquitectura actual de la Internet.
 \item Uno de los problemas que faltaría resolver es el direccionamiento IP, ya que, como ocurrió con la IPv4, la IPv6 quedará
 obsoleta en algún momento, y se necesita que la Internet del futuro resuelva dicho problema.
\end{itemize}




\begin{thebibliography}{1}

\bibitem{pursuitPDF}
``PURSUIT Publish Subscribe Internet Technology'', 2010.

\bibitem{UC}
University of Cambridge, ``Future internet aims to sever links with servers'', 2013.

\bibitem{Omi}
Adrian Raya, ``PURSUIT, el internet del futuro podría ser P2P'', 2013.

\bibitem{ps}
Lagutin, Dmitrij, Kari Visala, and Sasu Tarkoma. ``Publish/Subscribe for Internet: PSIRP Perspective.'' Future internet assembly 84 (2010).

\end{thebibliography}


\end{document}

